% Terminal historical appendix of the expansive study: the source home of
% the chronology that the concise companion's second page reprints. Dates,
% relation labels, alternatives and unresolved states are the generated
% projection of research/chronology.toml; the research entrypoint's preamble
% imports the generated definitions, once, so this file carries no import.
\section{Scriptural Date and Location}
\label{sec:date-location}
{\small Each passage the Mass reads or sings from Scripture, once, in
canonical order; beneath each, its attribution, its setting and the critical
horizon. The three orations are composed liturgical texts and carry no
biblical date; the Entrance is entered through its identified biblical basis,
and the acclamation through its adaptation of a biblical sentence.\par}

\begin{dossiertable}
Entrance (its basis) & Ps 37 (36):39--40 & No place named in the
verses & \chronodate{entrance}{\chronologyannotation{entrance}}\\
\cmidrule(lr){1-4}
\dossierprose{The composed antiphon \latin{Salus populi ego sum} is treated
through its identified biblical basis, Psalm 37:39--40, rather than as a
verbatim extract or a historical event. These verses specify no author,
place or occasion. No traditional date is held
for them. The NABRE Psalms introduction supplies the common pre-Maccabean
composition bound, not a precise date for this psalm or the Psalter's assembly.}
\midrule
Communion, Psalm option & Ps 119 (118):4--5 & No place named in the
verses & \chronodate{communion-ps119}{\chronologyannotation{communion-ps119}}\\
\cmidrule(lr){1-4}
\dossierprose{The opening petition of the great alphabetical psalm of the law:
a command received, followed by the wish that one's ways be directed to keep
it. No author is named and no traditional date is held; the date shown is the
Psalter's common critical composition bound (NABRE, Psalms introduction).
This is one of the Missal's two alternative
Communion texts for the day.}
\midrule
Responsorial Psalm & Ps 145 (144):2--3, 8--9, 17--18; response
v.~18a & No place named &
\chronodate{responsorial-psalm}{\chronologyannotation{responsorial-psalm}}\\
\cmidrule(lr){1-4}
\dossierprose{An alphabetical hymn to ``God my king'', headed in the
Douay--Rheims ``Praise, for David himself'' (v.~1): an attribution by title,
with no traditional date attached. The response is taken from v.~18a. The date
shown is the common composition bound in the NABRE Psalms introduction, which
does not securely date an individual psalm. Neither the title nor these
verses locates the writing.}
\midrule
First Reading & Isa 55:6--9 & Traditional ministry: Juda and Jerusalem (1:1);
critical horizon: Babylonian exile; writing place unestablished &
\chronodate{first-reading}{\chronologyannotation{first-reading}}\\
\cmidrule(lr){1-4}
\dossierprose{The two printed relations answer different questions and are not
to be merged. The first is the inherited attribution of the book to the
prophet Isaias: the range shown, 740--701~B.C., bounds his ministry in the
\work{Catholic Encyclopedia}'s traditional account (Souvay, 1910, ``Life''),
from the closing year of Ozias to the apparent latest prophecy, not his death.
It does not date this chapter's writing; Souvay's tentative collection
hypothesis concerns chapters 1--35. The second is the critical horizon of
the prophetic activity: the official NABRE introduction generally attributes chapters
40--55 to an anonymous poet prophesying toward the end of the Babylonian
exile. The prophecy addresses restoration from exile (Souvay, ``Second
Isaias''); that horizon is not a date for the writing or final assembly of
the book, nor proof of a particular writing site.}
\midrule
Gospel & Mt 20:1--16a & Spoken to the disciples on the journey toward
Jerusalem, after Peter's question (19:27) &
\chronodate{gospel}{\chronologyannotation{gospel}}\\
\cmidrule(lr){1-4}
\dossierevent{Narrated event: a parable of Jesus, told within the journey that
ends at Jerusalem, before the third prediction of the Passion (20:17--19); the
chronology corpus holds no date for the narrated event.}
\cmidrule(lr){1-4}
\dossierprose{Traditional attribution: St~Matthew the Apostle. The printed
figures are the disputed alternatives the \work{Catholic Encyclopedia}
reports, and they are not of equal standing. About A.D.~38--45 comes from early
ecclesiastical writers counting from the Ascension.
About A.D.~40--42 rests on a tradition of the Apostles' dispersal,
which the 1911 article itself calls ``admittedly not too reliable''.
A.D.~40--45 is the range reported by its survey of the Catholic critics of
its day. About A.D.~60--68 is a conditional reckoning that
fixes the Apostles' definitive departure later. About A.D.~64--67 rests
on a text of Irenaeus whose difficulties of interpretation the article says
prevent any positive conclusion. About A.D.~50 is Durand's date (1912) for the
Aramaic original only; of the Greek rendering, he writes, we know nothing
definite. The NABRE introduction states the modern critical horizon concisely:
the ``post-A.D.~70 date'', ``probably at least a decade later'', within its
probabilistic account; no closing year is given. Jacquier's ``Destination''
reports a Jewish-Christian audience and publication in Judea; the NABRE
instead regards the author as unknown and Antioch in Syria as a plausible
composition site for a community with Jewish and Gentile Christians. Neither
writing location is the setting of Jesus' parable.}
\midrule
Communion, John option & Jn 10:14 & Spoken to the Pharisees after the healing
of the man born blind (9:40--10:6) &
\chronodate{communion-jn10}{\chronologyannotation{communion-jn10}}\\
\cmidrule(lr){1-4}
\dossierevent{Narrated event: the discourse of the good shepherd; no event
date is held.}
\cmidrule(lr){1-4}
\dossierprose{Traditional attribution: St~John the Apostle. The two printed
ranges are the received disputed alternatives reported by the \work{Catholic
Encyclopedia}: Durand's approximate range for the Johannine writings (1912)
and Fonck's ``general opinion'', preceded by his admission that certain
historical dating is unavailable (1910, ``Circumstances of the composition'').
Fonck reports Ephesus as the writing place and Christians in Asia Minor as
the first readers; those are not the scene of Jesus' discourse.
The Missal adapts the verse. Christ's knowledge of his sheep belongs
to a discourse of self-offering, and this option stands as an alternative to
the Psalm text in the one Communion slot, not as a rite consecutive to it.}
\midrule
Acclamation (its basis) & cf.\ Acts 16:14b, an adaptation & Event: Philippi in
Macedonia, by the riverside on the sabbath (16:12--13) &
\chronodate{acclamation}{\chronologyannotation{acclamation}}\\
\cmidrule(lr){1-4}
\dossierevent{Narrated event: the conversion of Lydia during St~Paul's second
missionary journey (Acts 15:36--18:22).}
\cmidrule(lr){1-4}
\dossierprose{Traditional attribution: St~Luke. The two printed event ranges
are disputed alternatives the \work{Catholic Encyclopedia} gives for the
second missionary journey in its articles on St~Paul (1911) and on biblical
chronology (1908); the composition figure is the date its article on Acts
(1907) calls the most probable for the completion of the book. The mission's
date and the book's composition are distinct relations. Breen's probable
composition setting is Paul's Roman captivity, not Lydia's Philippi; Acts
addresses Theophilus (1:1), whose location is not supplied there. The liturgical verse
adapts Luke's narrative sentence into a petition for receptive hearing; it is
not itself a verse of Scripture.}
\midrule
Second Reading & Phil 1:20c--24, 27a & Paul in bonds (1:7, 13), at Rome in the
received chronology; to ``all the saints in Christ Jesus who are at Philippi''
(1:1) &
\chronodate{second-reading}{\chronologyannotation{second-reading}}\\
\cmidrule(lr){1-4}
\dossierprose{Traditional attribution: St~Paul, writing from imprisonment. The
printed figures are disputed alternatives within the Roman captivity: the
first is the year the \work{Catholic Encyclopedia}'s chronology of St~Paul
assigns the captivity letters (1911); the second, the range its article on the
letter gives for Paul ``at Rome'' (1911), where critics are said to disagree
whether the letter falls early or late in the Roman sojourn. The apostle's
present imprisonment frames his alternatives of continued service and being
with Christ; his hoped-for departure is not a dated narrated event.}
\end{dossiertable}
